\documentclass[12pt,a4paper]{book}
\usepackage[utf8]{inputenc}
\usepackage[T1]{fontenc}
\usepackage[french]{babel}
\usepackage{graphicx}
\usepackage{booktabs}
\usepackage{amsmath}
\usepackage{amssymb}
\usepackage{color}
\usepackage{hyperref}
\usepackage{geometry}
\usepackage{float}
\usepackage{caption}
\usepackage{subcaption}
\usepackage{fancyhdr}
\usepackage{setspace}
\usepackage{titlesec}
\usepackage{chngcntr}

% --- SUPPRESSION DES PAGES VIDES ---
\let\cleardoublepage\clearpage

% Configuration de la page
\geometry{
	a4paper,
	left=2.5cm,
	right=2.5cm,
	top=2.5cm,
	bottom=2.5cm
}

% Configuration des en-têtes et pieds de page
\pagestyle{fancy}
\fancyhf{}
\fancyhead[LE,RO]{\leftmark}
\fancyfoot[CE,CO]{\thepage}
\renewcommand{\headrulewidth}{0.4pt}

% Configuration des titres de chapitres
\titleformat{\chapter}[hang]{\Huge\bfseries}{}{0pt}{}
\titlespacing*{\chapter}{0pt}{-20pt}{20pt}

% Numérotation des figures par chapitre
\counterwithin{figure}{chapter}
\counterwithin{table}{chapter}
%emojie 
\usepackage{pifont}
\usepackage{xcolor}

% Informations du livre
\title{
	\begin{center}
		
\includegraphics[width=1.1in,height=1.1in]{./logo/Logo_fac_science}
		\hspace{9cm}
		\includegraphics[width=1.1in,height=1.1in]{./logo/logo_MIT}
	\end{center}
	\vspace{-2.5cm}
	\textbf{\textit{
			Université d'Antananarivo\\[12pt]
			Faculté des Sciences\\[15pt] 
			Mention Informatique et Technologie (MIT)}\\[50pt]} \bigskip
	\textsc{ grade	: \textbf{Master 1 }\\[2pt]}
	Parcours: \textsc{\textbf{ Innovation et Technologie (INT)}} 
	
	\vfill	
	\rule{0.75\textwidth}{3pt} \vspace{10pt}
	\textbf{\\ANALYSE EN COMPOSANTES PRINCIPALES SUR R :\\
	\Large \textbf{Étude des Habitudes de Vie Quotidienne}}
	\rule{0.75\textwidth}{3pt}
	\vfill
	
}

\date{\today \\[25pt]
	\textit{\textbf{Année universitaire : 2025 - 2026}}
}
\author{\underline{Nom:}\\[4pt]
	 \textbf{RABESIHARISON Tolotraniaina Henri}\\[1.5pt] 
}
\begin{document}
	
	% --- TABLE DES MATIÈRES ---
\pagenumbering{gobble}
\maketitle
\cleardoublepage
\pagenumbering{roman}
\addcontentsline{toc}{section}{Table des matières}
\tableofcontents
\cleardoublepage

\pagenumbering{arabic}
	% --- PRÉFACE (RACCOURCIE) ---
	\chapter*{Préface}
	\addcontentsline{toc}{chapter}{Préface}
	\markboth{Préface}{Préface}
	
	\vspace{0.5cm}
	
	Ce livre propose une analyse statistique des habitudes de vie à travers l'\textbf{Analyse en Composantes Principales (ACP)}. 
	
	L'étude porte sur \textbf{15 individus} et \textbf{11 variables} : sommeil, écran, sport, eau, travail, stress, alimentation, café, humeur, social et qualité du sommeil.
	
	\vspace{0.3cm}
	
	L'objectif est de :
	\begin{itemize}
		\item Visualiser les relations entre les habitudes.
		\item Identifier les profils types d'individus.
		\item Comprendre les facteurs clés du bien-être.
	\end{itemize}
	
	\vspace{0.3cm}
	
	Ce document s'adresse aux étudiants et professionnels souhaitant découvrir l'ACP à travers un cas pratique.
	
	\vspace{0.3cm}
	
	\noindent \textit{Bonne lecture !}
	
	% --- CHAPITRE 1 : INTRODUCTION ---
	\chapter{Introduction}
	\markboth{Introduction}{Introduction}
	
	\section{Contexte de l'étude}
	
	Les habitudes de vie quotidienne ont un impact majeur sur la santé et le bien-être des individus. La sédentarité, le stress, une mauvaise alimentation ou un sommeil insuffisant sont autant de facteurs qui peuvent affecter négativement la qualité de vie.
	
	\vspace{0.5cm}
	
	Dans ce contexte, il est essentiel de comprendre :
	\begin{itemize}
		\item \textbf{Quelles habitudes sont liées entre elles ?}
		\item \textbf{Quels sont les facteurs les plus influents ?}
		\item \textbf{Existe-t-il des profils types d'individus ?}
	\end{itemize}
	
	\section{Présentation des données}
	
	L'étude porte sur \textbf{15 individus} et \textbf{11 variables}. Le tableau ci-dessous présente un aperçu des données collectées :
	
	\vspace{0.5cm}
	
	\begin{table}[H]
		\centering
		\caption{Présentation des données étudiées}
		\begin{tabular}{lcccccc}
			\toprule
			\textbf{Individu} & \textbf{Sommeil} & \textbf{Ecran} & \textbf{Sport} & \textbf{Eau} & \textbf{Travail} & \textbf{Stress} \\
			\midrule
			Sarindra & 8.5 & 3 & 60 & 2.5 & 6 & 3 \\
			Bob & 6.0 & 6 & 20 & 1.8 & 8 & 7 \\
			Charlie & 7.0 & 4.5 & 40 & 2.0 & 7.5 & 5 \\
			Diana & 8.0 & 2.5 & 70 & 3.0 & 6 & 2 \\
			Elinoh & 5.5 & 7 & 10 & 1.5 & 9 & 8 \\
			Franck & 7.5 & 5 & 45 & 2.2 & 7 & 6 \\
			Miranto & 9.0 & 2 & 90 & 3.5 & 5 & 1 \\
			Faniry & 6.5 & 6.5 & 15 & 1.6 & 8.5 & 8 \\
			Toavina & 7.0 & 4 & 30 & 2.0 & 7 & 4 \\
			Jack & 5.0 & 8 & 5 & 1.2 & 10 & 10 \\
			Katy & 8.0 & 3.5 & 55 & 2.8 & 6.5 & 3 \\
			Leonard & 6.5 & 5.5 & 35 & 2.0 & 8 & 6 \\
			Mirotsy & 8.5 & 2.5 & 80 & 3.2 & 5.5 & 2 \\
			Noah & 7.0 & 4 & 45 & 2.5 & 7.5 & 5 \\
			Olivia & 6.0 & 7.5 & 25 & 1.8 & 9.5 & 9 \\
			\bottomrule
		\end{tabular}
	\end{table}
	
	\vspace{0.5cm}
	
	\section{Objectifs de l'étude}
	
	Les objectifs principaux de cette analyse sont :
	
	\begin{enumerate}
		\item \textbf{Identifier les relations} entre les différentes habitudes de vie.
		\item \textbf{Déterminer les facteurs} qui différencient le plus les individus.
		\item \textbf{Classifier} les individus en profils types.
		\item \textbf{Évaluer la qualité} de l'analyse statistique.
	\end{enumerate}
	
	\section{Structure du rapport}
	
	Ce rapport est structuré en six chapitres :
	
	\begin{enumerate}
		\item \textbf{Introduction} : Présentation du contexte, des données et des objectifs.
		\item \textbf{Méthodologie} : Explication de l'Analyse en Composantes Principales.
		\item \textbf{Cercle des corrélations} : Analyse des relations entre variables.
		\item \textbf{Contribution des variables} : Importance de chaque facteur.
		\item \textbf{Nuage des individus} : Profils types des personnes étudiées.
		\item \textbf{Valeurs propres} : Qualité et fiabilité de l'analyse.
	\end{enumerate}
	
	% --- CHAPITRE 2 : MÉTHODOLOGIE ---
	\chapter{Méthodologie}
	\markboth{Méthodologie}{Méthodologie}
	
	\section{L'Analyse en Composantes Principales (ACP)}
	
	L'\textbf{Analyse en Composantes Principales} est une méthode statistique multidimensionnelle qui permet de :
	
	\begin{itemize}
		\item \textbf{Réduire} la dimensionnalité des données.
		\item \textbf{Visualiser} les relations entre variables.
		\item \textbf{Identifier} les structures sous-jacentes.
	\end{itemize}
	
	\subsection{Principe de l'ACP}
	
	L'ACP transforme les variables initiales (corrélées entre elles) en de nouvelles variables non corrélées appelées \textbf{composantes principales}. Ces composantes sont classées par ordre d'importance décroissante.
	
	\vspace{0.5cm}
	
	\begin{center}
		\fbox{
			\begin{minipage}{0.8\textwidth}
				\textbf{Formule mathématique :}
				
				\[
				\begin{cases}
					C_1 = a_{11}X_1 + a_{12}X_2 + \dots + a_{1p}X_p \\
					C_2 = a_{21}X_1 + a_{22}X_2 + \dots + a_{2p}X_p \\
					\vdots \\
					C_p = a_{p1}X_1 + a_{p2}X_2 + \dots + a_{pp}X_p
				\end{cases}
				\]
				
				Où les $C_i$ sont les composantes principales et les $a_{ij}$ sont les coefficients (poids) calculés à partir des données.
			\end{minipage}
		}
	\end{center}
	
	\subsection{Étapes de l'ACP}
	
	Le processus d'analyse se déroule en plusieurs étapes :
	
	\begin{enumerate}
		\item \textbf{Standardisation} des données pour éviter l'influence des unités.
		\item \textbf{Calcul de la matrice de corrélation}.
		\item \textbf{Extraction des valeurs propres et des vecteurs propres}.
		\item \textbf{Calcul des composantes principales}.
		\item \textbf{Interprétation des résultats}.
	\end{enumerate}
	
	\section{Description des variables}
	
	Le tableau ci-dessous présente les variables utilisées dans cette étude :
	
	\begin{table}[H]
		\centering
		\caption{Description des variables}
		\begin{tabular}{lcp{8cm}}
			\toprule
			\textbf{Variable} & \textbf{Unité} & \textbf{Description} \\
			\midrule
			Sommeil & Heures & Durée de sommeil par nuit \\
			Ecran & Heures & Temps passé devant un écran \\
			Sport & Minutes & Temps de pratique sportive quotidienne \\
			Eau & Litres & Consommation d'eau par jour \\
			Travail & Heures & Temps de travail par jour \\
			Stress & /10 & Niveau de stress ressenti \\
			Alimentation & /10 & Qualité de l'alimentation \\
			Café & Tasses & Consommation de café par jour \\
			Humeur & /10 & Niveau de bien-être général \\
			Social & Heures & Temps d'interactions sociales \\
			Sommeil\_qualite & /10 & Qualité subjective du sommeil \\
			\bottomrule
		\end{tabular}
	\end{table}
	
	\section{Justification de l'ACP pour cette étude}
	
	L'ACP est particulièrement adaptée à cette étude car :
	
	\begin{itemize}
		\item Les variables sont \textbf{nombreuses} (11 variables).
		\item Il existe probablement des \textbf{corrélations} entre les variables.
		\item Nous souhaitons \textbf{visualiser} les relations de manière simple.
		\item Nous voulons \textbf{identifier} des profils types d'individus.
	\end{itemize}
	
	% --- CHAPITRE 3 : CERCLE DES CORRÉLATIONS ---
	\chapter{Cercle des Corrélations}
	\markboth{Cercle des Corrélations}{Cercle des Corrélations}
	
	\section{Présentation du graphique}
	
	Le cercle des corrélations est un outil visuel essentiel de l'ACP. Il permet de visualiser :
	
	\begin{itemize}
		\item Les \textbf{corrélations} entre les variables.
		\item La \textbf{contribution} de chaque variable aux axes principaux.
	\end{itemize}
	
	\begin{figure}[H]
		\centering
		\includegraphics[width=0.8\textwidth]{./img/cercle_corelation.png}
		\caption{Cercle des corrélations - Relations entre les habitudes de vie}
		\label{fig:cercle}
	\end{figure}
	
	\section{Interprétation des axes}
	
	\subsection{Axe horizontal (Dimension 1)}
	
	L'axe horizontal est l'axe principal qui explique \textbf{52\%} de la variance totale.
	
	\begin{itemize}
		\item \textbf{À droite} de l'axe : Variables associées à un style de vie "stressant et sédentaire" :
		\begin{itemize}
			\item \textcolor{red}{Travail}
			\item \textcolor{red}{Ecran}
			\item \textcolor{red}{Stress}
		\end{itemize}
		
		\item \textbf{À gauche} de l'axe : Variables associées à un style de vie "sain et équilibré" :
		\begin{itemize}
			\item \textcolor{blue}{Sport}
			\item \textcolor{blue}{Alimentation}
			\item \textcolor{blue}{Sommeil\_qualite}
			\item \textcolor{blue}{Humeur}
		\end{itemize}
	\end{itemize}
	
	\subsection{Axe vertical (Dimension 2)}
	
	L'axe vertical explique \textbf{21\%} de la variance totale. Il permet d'affiner les profils :
	
	\begin{itemize}
		\item Il sépare les personnes qui dorment bien mais sont stressées de celles qui ne dorment pas bien mais sont calmes.
		\item Il peut être lié à des facteurs comme la \textbf{consommation de café} ou les \textbf{interactions sociales}.
	\end{itemize}
	
	\section{Groupes de variables identifiés}
	
	\begin{table}[H]
		\centering
		\caption{Groupes de variables identifiés}
		\begin{tabular}{p{5cm}p{5cm}}
			\toprule
			\textbf{Groupe "Stress/Sédentarité"} & \textbf{Groupe "Bien-être"} \\
			\midrule
			\textcolor{red}{Travail} & \textcolor{blue}{Sport} \\
			\textcolor{red}{Ecran} & \textcolor{blue}{Alimentation} \\
			\textcolor{red}{Stress} & \textcolor{blue}{Sommeil\_qualite} \\
			\textcolor{red}{Café} & \textcolor{blue}{Humeur} \\
			& \textcolor{blue}{Sommeil} \\
			& \textcolor{blue}{Eau} \\
			\bottomrule
		\end{tabular}
	\end{table}
	
	\section{Conclusion du chapitre}
	
	Le cercle des corrélations met clairement en évidence \textbf{deux groupes opposés} d'habitudes de vie :
	
	\begin{enumerate}
		\item Un groupe \textbf{négatif} associé au stress, au travail et à l'écran.
		\item Un groupe \textbf{positif} associé au sport, à une bonne alimentation et à un sommeil de qualité.
	\end{enumerate}
	
	Ces deux groupes sont fortement opposés le long de l'axe horizontal (Dimension 1).
	
	% --- CHAPITRE 4 : CONTRIBUTION DES VARIABLES ---
	\chapter{Contribution des Variables}
	\markboth{Contribution des Variables}{Contribution des Variables}
	
	\section{Présentation du graphique}
	
	Ce graphique montre la \textbf{contribution de chaque variable} à la construction des axes principaux.
	
	\begin{figure}[H]
		\centering
		\includegraphics[width=0.8\textwidth]{./img/contribution.png}
		\caption{Contribution des variables à l'analyse}
		\label{fig:contribution}
	\end{figure}
	
	\section{Analyse des résultats}
	
	\subsection{Répartition des contributions}
	
	Le tableau ci-dessous présente les contributions détaillées de chaque variable :
	
	\begin{table}[H]
		\centering
		\caption{Contributions des variables}
		\begin{tabular}{lcc}
			\toprule
			\textbf{Variable} & \textbf{Contribution (\%)} & \textbf{Rang} \\
			\midrule
			Travail & 16.9\% & 1 \\
			Ecran & 16.8\% & 2 \\
			Stress & 16.8\% & 3 \\
			Sommeil & 16.6\% & 4 \\
			Sport & 16.5\% & 5 \\
			Eau & 16.4\% & 6 \\
			\bottomrule
		\end{tabular}
	\end{table}
	
	\subsection{Interprétation}
	
	\begin{itemize}
		\item Toutes les variables contribuent de manière \textbf{quasi-équilibrée} (environ 16-17\%).
		\item Les 3 variables les plus influentes sont :
		\begin{enumerate}
			\item \textbf{Travail} (16.9\%)
			\item \textbf{Ecran} (16.8\%)
			\item \textbf{Stress} (16.8\%)
		\end{enumerate}
		\item Le \textbf{Sport} (16.5\%) et le \textbf{Sommeil} (16.6\%) sont également des facteurs clés.
		\item La \textbf{consommation d'eau} (16.4\%) ferme la marche des variables principales.
	\end{itemize}
	
	\section{Conclusion du chapitre}
	
	\begin{center}
		\fbox{
			\begin{minipage}{0.85\textwidth}
				\centering
				\textbf{La qualité de vie est un phénomène multifactoriel.}
				\vspace{0.2cm}
				
				Aucun facteur ne domine complètement les autres. Toutes les variables ont une importance similaire dans l'explication des différences entre les individus.
			\end{minipage}
		}
	\end{center}
	
	\vspace{0.5cm}
	
	Cette répartition équilibrée montre que pour améliorer son bien-être, il est nécessaire d'agir simultanément sur \textbf{plusieurs aspects} de la vie quotidienne : le travail, le stress, le sommeil, le sport et l'alimentation.
	
	% --- CHAPITRE 5 : NUAGE DES INDIVIDUS ---
	\chapter{Nuage des Individus}
	\markboth{Nuage des Individus}{Nuage des Individus}
	
	\section{Présentation du graphique}
	
	Le nuage de points permet de \textbf{visualiser la position de chaque individu} dans l'espace défini par les axes principaux.
	
	\begin{figure}[H]
		\centering
		\includegraphics[width=0.8\textwidth]{./img/nuage_de_point.png}
		\caption{Nuage de points - Taille selon le Stress, Couleur selon le Sport}
		\label{fig:nuage}
	\end{figure}
	
	\section{Légende du graphique}
	
	\begin{itemize}
		\item \textbf{Couleur des points} : Indique le niveau de pratique sportive
		\begin{itemize}
			\item \textcolor{blue}{Bleu} = Beaucoup de sport
			\item \textcolor{red}{Rouge} = Peu de sport
		\end{itemize}
		\item \textbf{Taille des points} : Indique le niveau de stress
		\begin{itemize}
			\item \textcolor{red}{\Large •} = Stress élevé
			\item \textcolor{blue}{\small •} = Stress faible
		\end{itemize}
	\end{itemize}
	
	\section{Identification des profils types}
	
	\subsection{Profil 1 : Le mode de vie "sain" (Gauche du graphique)}
	
	Ces personnes se caractérisent par :
	
	\begin{itemize}
		\item \textcolor{blue}{Beaucoup de sport} (couleur bleue)
		\item \textcolor{blue}{Faible stress} (petite taille)
		\item Moins de temps d'écran
		\item Meilleure alimentation
		\item Bonne qualité de sommeil
	\end{itemize}
	
	\textbf{Exemples d'individus :}
	\begin{enumerate}
		\item \textbf{Miranto} : 90 min de sport, stress 1/10
		\item \textbf{Mirotsy} : 80 min de sport, stress 2/10
		\item \textbf{Diana} : 70 min de sport, stress 2/10
	\end{enumerate}
	
	\subsection{Profil 2 : Le mode de vie "stressé" (Droite du graphique)}
	
	Ces personnes se caractérisent par :
	
	\begin{itemize}
		\item \textcolor{red}{Peu de sport} (couleur rouge)
		\item \textcolor{red}{Stress élevé} (grande taille)
		\item Beaucoup de temps d'écran
		\item Forte charge de travail
		\item Moins bonne qualité de sommeil
	\end{itemize}
	
	\textbf{Exemples d'individus :}
	\begin{enumerate}
		\item \textbf{Jack} : 5 min de sport, stress 10/10
		\item \textbf{Olivia} : 25 min de sport, stress 9/10
		\item \textbf{Elinoh} : 10 min de sport, stress 8/10
	\end{enumerate}
	
	\section{Analyse des cas extrêmes}
	
	\begin{table}[H]
		\centering
		\caption{Cas extrêmes de l'étude}
		\begin{tabular}{lcccc}
			\toprule
			\textbf{Individu} & \textbf{Ecran} & \textbf{Sport} & \textbf{Stress} & \textbf{Profil} \\
			\midrule
			\textcolor{red}{Jack} & 8h & 5 min & 10/10 & \textcolor{red}{Stressé, sédentaire} \\
			\textcolor{red}{Olivia} & 7.5h & 25 min & 9/10 & \textcolor{red}{Stressée, sédentaire} \\
			\textcolor{red}{Elinoh} & 7h & 10 min & 8/10 & \textcolor{red}{Stressée, sédentaire} \\
			\midrule
			\textcolor{blue}{Miranto} & 2h & 90 min & 1/10 & \textcolor{blue}{Sain, sportif} \\
			\textcolor{blue}{Mirotsy} & 2.5h & 80 min & 2/10 & \textcolor{blue}{Sain, sportif} \\
			\textcolor{blue}{Diana} & 2.5h & 70 min & 2/10 & \textcolor{blue}{Sain, sportif} \\
			\bottomrule
		\end{tabular}
	\end{table}
	
	\section{Conclusion du chapitre}
	
	Le nuage de points confirme la \textbf{séparation nette} entre deux styles de vie :
	
	\begin{enumerate}
		\item \textbf{Style de vie sain} : Sportif, faible stress, bien-être.
		\item \textbf{Style de vie stressé} : Sédentaire, stress élevé, mauvaise qualité de vie.
	\end{enumerate}
	
	Les cas extrêmes (Jack, Olivia, Miranto, Mirotsy) sont les principaux "moteurs" de l'analyse. Ils créent les deux groupes distincts observés.
	
	% --- CHAPITRE 6 : VALEURS PROPRES ---
	\chapter{Valeurs Propres - Qualité de l'Analyse}
	\markboth{Valeurs Propres}{Valeurs Propres}
	
	\section{Présentation du graphique}
	
	Les valeurs propres indiquent \textbf{l'importance de chaque axe} (dimension) dans l'analyse.
	
	\begin{figure}[H]
		\centering
		\includegraphics[width=0.8\textwidth]{./img/valeurs_propres.png}
		\caption{Valeurs propres - Pourcentage de variance expliquée}
		\label{fig:vp}
	\end{figure}
	
	\section{Définition des valeurs propres}
	
	Les \textbf{valeurs propres} (ou "eigenvalues") sont des nombres qui mesurent l'importance de chaque axe.
	
	\begin{itemize}
		\item \textbf{Plus la valeur propre est grande} = Plus l'axe est important.
		\item La somme de toutes les valeurs propres = 100\% de l'information.
	\end{itemize}
	
	\begin{center}
		\fbox{
			\begin{minipage}{0.8\textwidth}
				\centering
				\textbf{Formule :}
				\[
				\lambda_i = \text{Variance expliquée par la dimension } i
				\]
				\[
				\text{Pourcentage expliqué} = \frac{\lambda_i}{\sum \lambda_j} \times 100
				\]
			\end{minipage}
		}
	\end{center}
	
	\section{Résultats de l'analyse}
	
	\subsection{Pourcentages de variance expliquée}
	
	\begin{table}[H]
		\centering
		\caption{Résultats des valeurs propres}
		\begin{tabular}{ccc}
			\toprule
			\textbf{Dimension} & \textbf{Variance expliquée (\%)} & \textbf{Cumul (\%)} \\
			\midrule
			Dimension 1 & 52.0\% & 52.0\% \\
			Dimension 2 & 21.0\% & 73.0\% \\
			Dimension 3 & 1.2\% & 74.2\% \\
			Dimension 4 & 0.4\% & 74.6\% \\
			Dimension 5 & 0.3\% & 74.9\% \\
			\bottomrule
		\end{tabular}
	\end{table}
	
	\subsection{Interprétation des résultats}
	
	\begin{enumerate}
		\item \textbf{Dimension 1} (52\%)
		\begin{itemize}
			\item C'est l'axe le plus important.
			\item Il oppose les deux styles de vie : "bien-être" vs "stress/sédentarité".
			\item C'est l'axe horizontal de tous nos graphiques.
		\end{itemize}
		
		\item \textbf{Dimension 2} (21\%)
		\begin{itemize}
			\item C'est le deuxième axe le plus important.
			\item Il affine les profils des individus.
			\item C'est l'axe vertical de nos graphiques.
		\end{itemize}
		
		\item \textbf{Dimensions 3, 4, 5} (< 1\%)
		\begin{itemize}
			\item Ces axes apportent très peu d'information.
			\item Ils peuvent être ignorés sans perte significative.
		\end{itemize}
	\end{enumerate}
	
	\section{Utilité des valeurs propres}
	
	\subsection{1. Évaluer la qualité de l'analyse}
	
	\begin{center}
		\fbox{
			\begin{minipage}{0.85\textwidth}
				\centering
				\textbf{Bonne qualité de l'analyse}
				\vspace{0.2cm}
				
				Les 2 premiers axes résument \textbf{73\%} des données.
				\vspace{0.2cm}
				
				C'est un excellent résultat qui valide la fiabilité de notre analyse.
			\end{minipage}
		}
	\end{center}
	
	\subsection{2. Choisir le nombre d'axes à garder}
	
	Règle pratique : on garde les dimensions dont la valeur propre est :
	\begin{itemize}
		\item \textbf{Supérieure à 1} (critère de Kaiser).
		\item Ou qui cumulent environ \textbf{70-80\%} de la variance.
	\end{itemize}
	
	Ici, on garde :
	\begin{itemize}
		\item \textbf{Dimension 1} : Valeur propre > 1
		\item \textbf{Dimension 2} : Valeur propre > 1
		\item Les autres : Valeur propre < 1 → On les ignore.
	\end{itemize}
	
	\subsection{3. Valider les représentations graphiques}
	
	\begin{itemize}
		\item Le cercle des corrélations est projeté sur les axes 1 et 2.
		\item Ces axes expliquent 73\% de l'information.
		\item Donc les graphiques sont \textbf{fiables} et \textbf{représentatifs}.
	\end{itemize}
	
	\section{Conclusion du chapitre}
	
	\begin{center}
		\fbox{
			\begin{minipage}{0.85\textwidth}
				\centering
				\textbf{Points clés à retenir :}
				\vspace{0.2cm}
				
				\begin{enumerate}
					\item La Dimension 1 (52\%) oppose "bien-être" et "stress".
					\item La Dimension 2 (21\%) affine les profils.
					\item Les 2 axes résument 73\% de l'information → analyse fiable.
					\item Les autres dimensions (< 1\%) sont négligeables.
				\end{enumerate}
			\end{minipage}
		}
	\end{center}
	
	% --- CONCLUSION GÉNÉRALE (SIMPLIFIÉE) ---
	\chapter*{Conclusion}
	\addcontentsline{toc}{chapter}{Conclusion}
	\markboth{Conclusion}{Conclusion}
	
	\vspace{0.5cm}
	
	\begin{center}
		\Large \textbf{En résumé}
	\end{center}
	
	\vspace{0.3cm}
	
	L'ACP a permis d'identifier \textbf{deux profils} de vie :
	
	\vspace{0.3cm}
	% Définition de commandes personnalisées
	\newcommand{\ok}{\textcolor{green}{\ding{51}}}     % ✅ vert
	\newcommand{\ko}{\textcolor{red}{\ding{55}}}        % ❌ rouge
	\newcommand{\warning}{\textcolor{orange}{\ding{43}}} % ⚠️ orange
	\newcommand{\info}{\textcolor{blue}{\ding{42}}} 
	¨
	\begin{table}[H]
		\centering
		\begin{tabular}{p{5cm}p{5cm}}
			\toprule
			\textcolor{blue}{\textbf{Profil Sain}} & \textcolor{red}{\textbf{Profil Stressé}} \\
			\midrule
			\ok Beaucoup de sport & \ko Peu de sport \\
			\ok Faible stress & \ko Stress élevé \\
			\ok Peu d'écran & \ko Beaucoup d'écran \\
			\ok Bonne alimentation & \ko Mauvaise alimentation \\
			\ok Sommeil de qualité & \ko Sommeil médiocre \\
			\bottomrule
		\end{tabular}
	\end{table}
	
	\vspace{0.3cm}
	
	\textbf{Les principaux enseignements :}
	
	\begin{itemize}
		\item Le \textcolor{red}{stress, le travail et l'écran} sont les facteurs les plus négatifs.
		\item Le \textcolor{blue}{sport et l'alimentation} sont les clés du bien-être.
		\item L'analyse est \textbf{très fiable} (73\% d'information résumée).
	\end{itemize}
	
	\vspace{0.3cm}
	
	\begin{center}
		\fbox{
			\begin{minipage}{0.8\textwidth}
				\centering
				\Large \textbf{Moins d'écran + Plus de sport = Meilleure santé}
			\end{minipage}
		}
	\end{center}
	
	\vspace{0.5cm}
	\begin{flushright}
		\textit{Juillet 2026}
	\end{flushright}
	
	% --- INDEX ---
	\chapter*{Index}
	\addcontentsline{toc}{chapter}{Index}
	\markboth{Index}{Index}
	
	\begin{description}
		\item[ACP] Analyse en Composantes Principales
		\item[Corrélations] Relations entre les variables
		\item[Dimension 1] Axe horizontal (52\% de variance)
		\item[Dimension 2] Axe vertical (21\% de variance)
		\item[Nuage de points] Représentation des individus
		\item[Profil sain] Individus sportifs, faible stress
		\item[Profil stressé] Individus sédentaires, stress élevé
		\item[Valeurs propres] Importance des axes
		\item[Variance] Information contenue dans les données
	\end{description}
	
\end{document}